% Owned proposed B.140. No checkpoint28 source or review is altered.
\subsection*{Actual retained-section corrections and an obstruction to every inner repair}

\statement{B.140}{Proved}{The two named retained sections fail to be source-invariant after the whole filling, and no family-kernel adjustment of the outer source repairs either section, discretely or in full completion}
\label{stmt:B140}
Use the actual B.123 groups
\[
 \begin{gathered}
 F=F(U_1,\ldots,U_8)*T,\quad T=F(f,g,x),\quad L=\ker(F\to T),\\
 W=F/N,\qquad J=\ker(W\to T),\qquad \Gamma_0=W\rtimes S.
 \end{gathered}
\]
where $S=F(m,v)$ and $N$ is the entire forty-seed normal closure
under all signed source and retained-deck translates.
The family labels, in the original $F_{11}$ basis, are
\[
 (U_1,\ldots,U_8)=(a,B,d,e,\bar a,\bar B,\bar d,\bar e),
 \qquad \bar u=cuc,\quad g=cfc,\quad x=\tau c.
\]
By B.127, $N\subseteq[L,L]$ and
$J_{\rm ab}=M=\mathbb Z[T]^8$ with its left deck action.

Consider the following two actual sections of $W\to T$:
\begin{equation}\label{eq:B140-sections}
 \begin{aligned}
 \sigma_o(f,g,x)&=(f,g,x),\\
 \sigma_\ell(f,g,x)
   &=(f,a^{-1}\bar a g\bar a^{-1}a,\ x\bar a^{-1}a)
     =(f,CfC,\tau C),\qquad C=a^{-1}ca .
 \end{aligned}
\end{equation}
Both named subgroups, and their conjugates in $\Gamma_0$,
have their full induced profinite topology.
The old section is not preserved by either signed $m$ or
either signed $v$ action. The local section is preserved
by both signed $m$ actions but not by either signed $v$ action.
All these noninvariance assertions hold after the entire
node filling, and after full completion.

More strongly, there is no $j\in J$ for which $jv$
normalizes either subgroup in \eqref{eq:B140-sections}.
There is likewise no $j\in\overline J$ for which $jv$
normalizes either actual completed subgroup.
These exclusions also apply after replacing a named section
by any common $J$-conjugate, or completed section by any
common $\overline J$-conjugate.
In particular a section of the whole coarse quotient cannot
restrict to either named retained section or these conjugates.
An arbitrary nonconjugate retained section is not excluded.
No original fixed-source equality is decided.

\begin{proof}
\emph{B.140.1. Actual post-node defect classes.}
Let $\phi_s$ denote the automorphism induced on $T$ by an
actual signed source $s$. For either section $\sigma$, put
\begin{equation}\label{eq:B140-defect}
 D_{\sigma,s}(t)
   =\alpha_s(\sigma(t))\sigma(\phi_s(t))^{-1}\in J,\qquad
 d_{\sigma,s}(t)=[D_{\sigma,s}(t)]\in M .
\end{equation}
These are classes in the actual post-node kernel, because
$N\subseteq[L,L]$ holds for the whole defining normal closure.
Conjugation by $J$ is trivial on $M$; hence either section
induces the same left regular $T$-action.

The coefficient extraction is completely specified by the
free-product Schreier scan of B.127. Read a word in $F$,
retaining its current reduced $T$ prefix. A family letter
$U_i^\epsilon$ emits $\epsilon e_{(t,i)}$ at that prefix,
and a retained letter updates it. The finite sum of emitted
terms is its exact class in $M$, once its retained head is
trivial. Apply this to each word in \eqref{eq:B140-defect},
using the complete signed operators and both actual section
substitutions. In the $F_{11}$ numbering the local section is
\[
 (5),\quad(-1,6,10,-6,1),\quad(11,-6,1).
\]
The complete input check verifies fourteen inverse regional
basis substitutions, all forty-four signed generator
comparisons between the two independently marked operator
tables, and twelve complete earlier retained Schreier rewrites.
It then extracts all twenty-four defect classes.
Their support counts, in the order $f,g,x$, are
\begin{equation}\label{eq:B140-support-table}
 \begin{array}{c|c|c}
 s&\sigma_o&\sigma_\ell\\ \hline
 m&(4,4,4)&(0,0,0)\\
 m^{-1}&(4,6,4)&(0,0,0)\\
 v&(12,28,12)&(12,28,12)\\
 v^{-1}&(12,24,12)&(12,24,12).
 \end{array}
\end{equation}
The zero entries for the local $m$ actions are also literal
equalities from B.135's partial-inner formulas, so invariance
there is not inferred from abelian vanishing.

For explicit nonzero witnesses let
$q=xgx^{-1}$ and write $e_a,e_{\bar a}$ for the indicated
family basis vectors. The first old-section defects are
\[
 d_{\sigma_o,m}(f)
     =f^{-1}(g-1)(e_a-e_{\bar a}),\qquad
 d_{\sigma_o,m^{-1}}(f)
     =(g-1)(e_a-e_{\bar a}).
\]
For either section, the entire $B$-color components satisfy
\begin{equation}\label{eq:B140-minimal-input}
 \begin{gathered}
 \phi_v(f)=q,\qquad
 h=\phi_v(g)=f^{-1}gx^{-1}fxg^{-1}f,\qquad
 \phi_v(x)=xg^{-1}f,\\
 (d_{\sigma,v}(f))_B=1-q^{-1},\qquad
 [1]\,(d_{\sigma,v}(g))_B=1,\qquad
 (d_{\sigma,v^{-1}}(f))_B=(q-1)f .
 \end{gathered}
\end{equation}
Here $[1]$ denotes the identity coefficient of the displayed
finite group-ring element. The complete first component and
the specified second coefficient were separately checked
directly from the actual extraction.
These nonzero classes already prove the discrete
noninvariance assertions.

\emph{B.140.2. A common inner correction cannot solve two generator equations.}
Suppose $jv$ normalized $\sigma(T)$.
Its induced action there must be $\phi_v$, since the filling
projection restricts to an isomorphism on that section.
Writing $\mu=[j]\in M$, the required equations in $M$ are
\begin{equation}\label{eq:B140-common-equations}
 d_{\sigma,v}(t)=(\phi_v(t)-1)\mu\qquad(t=f,g,x).
\end{equation}
Indeed inner adjustment changes the defect to
$\mu+d_{\sigma,v}(t)-\phi_v(t)\mu$.
The same $\mu$ must work in all three equations.

Left multiplication by $q-1$ is injective on $\mathbb Z[T]$:
a finitely supported vector fixed by the infinite-order
translation $q$ vanishes on every infinite cyclic orbit.
Since $(q-1)q^{-1}=1-q^{-1}$,
the first $B$-color equation uniquely forces
$\mu_B=q^{-1}$.
The second would require
\[
 (d_{\sigma,v}(g))_B=hq^{-1}-q^{-1}.
\]
Its right side has identity coefficient zero, because
$q\neq1$ and the displayed reduced words give $h\neq q$.
Its left side has identity coefficient one by
\eqref{eq:B140-minimal-input}. This contradiction excludes
every $j\in J$, without a restriction on its length or
selected-involution class.

For completeness, the whole three-equation check uses
finite cyclic-line division, not a search for correcting words.
For $q_s=\phi_s(f)$, pull support words back by $\phi_s^{-1}$
and write them uniquely as $f^nr$, where $r$ does not begin
in $f^{\pm1}$. On each such line, the equation is
$\mu(n-1)-\mu(n)=d(n)$.
A finite-support solution exists exactly when
$\sum_n d(n)=0$, and is then uniquely
$\mu(n)=-\sum_{k\leq n}d(k)$.
The actual first equation has a candidate in all eight
section/operator cases. The candidate has two terms for
old $m^{\pm1}$, zero for local $m^{\pm1}$, and six for
$v^{\pm1}$. Subtracting it from the other two equations
gives, respectively, support counts
$(6,4),(6,6),(0,0),(0,0),(36,16),(24,14)$,
with the last two identical for both sections.
The two-equation proof above needs only
\eqref{eq:B140-minimal-input}.

\emph{B.140.3. Actual finite targets and the topology of both sections.}
Set
\[
 P=F/[L,L]=M\rtimes T .
\]
This is a finitely generated actual quotient, with the eight
family and three retained generators. Both signed source
actions preserve $[L,L]$, and $N\subseteq[L,L]$.
Thus there is an actual quotient map
\begin{equation}\label{eq:B140-metabelian-quotient}
 \Gamma_0\longrightarrow P\rtimes S .
\end{equation}
It retains the nonlinear source corrections in its induced
automorphisms; it does not assert that $S$ preserves the
chosen $T$ section in $P$.

For $n\geq1$ and $U\triangleleft T$ of finite index, put
\[
 P_{n,U}=\bigl((\mathbb Z/n)[T/U]^8\bigr)\rtimes(T/U).
\]
These finite wreath heads are cofinal among all finite
targets of $P$. To see this, for $\rho:P\to A$ finite
choose $n$ divisible by the exponent of $\rho(M)$ and
$U\leq\ker(\rho|_T)$. The action on $\rho(M)$ factors
through $T/U$. Evaluation of the eight basis vectors
at $\rho(U_i)$ and retained letters at their $\rho$ images
defines a homomorphism $P_{n,U}\to\rho(P)$ factoring $\rho$.
This remains valid if $\rho(M)$ and $\rho(T)$ intersect.
It follows that
\begin{equation}\label{eq:B140-completed-P}
 \widehat P=\Omega\rtimes\widehat T,\qquad
 \Omega=\widehat{\mathbb Z}[[\widehat T]]^8 .
\end{equation}

Also $P$ has its full induced topology in $P\rtimes S$.
Given any prescribed finite-index kernel of $P$, choose
a characteristic finite-index subgroup inside it, using
finite generation. The induced $S$-action on the finite
quotient has finite image. A finite quotient of $S$
retaining that action gives the required finite
semidirect target. This supplies whole-group finite
refinements even when an initial wreath head itself is
not source-invariant.

Each named section maps to a section of $P\to T$.
An arbitrary finite target of its free $T$ therefore
extends to a finite target of $P$ by that retraction,
and then extends in the finite-kernel-containment sense
to $P\rtimes S$. Composing with
\eqref{eq:B140-metabelian-quotient} proves the full
induced topology of that actual section in $\Gamma_0$.
The same holds for its conjugates.
Thus the named completed subgroups are actual embedded
copies of $\widehat T$.

The map \eqref{eq:B140-metabelian-quotient} extends to full
completions. Its image on $\overline J$ lies in the
actual module factor $\Omega$ of \eqref{eq:B140-completed-P}.
Indeed the discrete $J$ maps into $M$, and the latter
has that closure in $\widehat P$, embedded in
$\widehat{P\rtimes S}$ by the preceding finite refinements.
Accordingly any completed normalizing lift must satisfy
the same necessary equations
\eqref{eq:B140-common-equations} in $\Omega$.
No assertion about $\widehat J$ or the continuous
abelianization of $\overline J$ is needed.

\emph{B.140.4. Primitive translations remain injective on the completed module.}
We prove that $q-1$ is injective on
$\widehat{\mathbb Z}[[\widehat T]]$.
The primitive element $q=xgx^{-1}$ has an integral
character $\chi_q:T\to\mathbb Z$ with $\chi_q(q)=1$.
Suppose $q\lambda=\lambda$.
Fix a finite quotient $F_0$ of $T$ and a coefficient
modulus $n$. If $q$ has order $d$ in $F_0$, let $F_1$
be the actual image of
\[
 T\longrightarrow F_0\times C_{dn},
 \qquad t\longmapsto(t,\chi_q(t)\bmod dn).
\]
The image of $q$ has order exactly $dn$ in $F_1$;
its $d$th power lies in the kernel of $F_1\to F_0$
and has order $n$.
Within each fiber, its left action has cycles of length
$n$. The image of $\lambda$ in $(\mathbb Z/n)[F_1]$
has constant coefficients on these cycles.
Pushing forward therefore gives zero in
$(\mathbb Z/n)[F_0]$.
This holds for every $F_0,n$, proving $\lambda=0$.
No group-ring domain assertion is used.

The first completed equation again forces
$\mu_B=q^{-1}$. The second has the same explicit
nonzero finite-support residual.
Such a residual stays nonzero in the completed group
ring: residual finiteness of free $T$ separates its
finitely many support words in one finite quotient,
and a coefficient modulus retains a nonzero coefficient.
This uses the identity coefficient only in a finite
integral group-ring element; it does not assume a
continuous point-coefficient functional on the
completed ring.
The completed contradiction follows. The same finite
target separation also promotes every nonzero defect
in \eqref{eq:B140-support-table} to actual completed
noninvariance through \eqref{eq:B140-metabelian-quotient}.

\emph{B.140.5. Common conjugates and the exact limitation.}
If $k\sigma(T)k^{-1}$, $k\in J$, were normalized by a
lift $jv$, then $k^{-1}(jv)k$ would be a lift of the
same coarse $v$ normalizing $\sigma(T)$.
This is impossible. The identical argument holds
for $j,k\in\overline J$ and the completed subgroups.
A normalizing lift of any coarse $tv$, $t\in T$,
would likewise yield a forbidden lift of $v$ after
multiplication by $\sigma(t)^{-1}$.
Thus the claimed section restrictions are excluded.
An arbitrary different retained section need not be
a common conjugate of either named section; it has
not been ruled out. These obstructions do not decide
either remaining fixed-kernel equality of B.139.
\end{proof}

The full actual correction extraction and the independent
exact inner-equation checks are retained with genuine outputs.
No source-orbit enumeration or finite-group census is used.
\begin{quote}\footnotesize\raggedright
\path{code/check_root_coarse_corrections29.py},
\path{records/root_coarse_corrections29.json},
\path{work_CE/code/check_section_inner_adjustment29.py},
\path{work_CE/next_symmetric_nodes/section_inner_adjustment29.json},
\path{work_CE/code/check_section_inner_minimal29.py},
\path{work_CE/retained_section_inner_obstruction29.txt}.
\end{quote}
